\section{K. The ten programme questions on 26/011}
\label{sec:geometry-K}

These are headings in a research programme, not ten fully quantified
conjectures. In particular, ``modules'' means moduli, not modules over a
ring. We retain the order and distinguish the actual heading from each
precise mathematical assertion that it suggests. The readings
``automorphisme(s)'', ``Hilbert'', and $\mathrm{GL}_n$ do not change the
arguments. The uncertain parenthesis in (f), ``par modules / par un autre'',
does affect how ambitious a classification is requested; both descent
classification and representability of the moduli stack are treated below.

\subsection{(a) Rings of definition for structures without automorphisms}

\begin{proposition}[Proved: uniqueness removes the cocycle obstruction]
\label{geo-Ka}
Let $\mathcal X$ be a stack for a topology in which descent is effective,
and let $U\to S$ be a covering. Suppose $\xi\in\mathcal X(U)$ has trivial
automorphism sheaf after every base change, and its two pullbacks to
$U\times_S U$ are locally isomorphic. Then $\xi$ descends to $S$, uniquely
up to a unique isomorphism compatible with the given object on $U$.
In particular, an object over a finite Galois extension $L/k$ whose Galois
conjugates are isomorphic and whose geometric automorphism group is trivial
has a unique effective descent to $k$ in such a stack.
\end{proposition}
\begin{proof}
Between two locally isomorphic objects with trivial automorphism sheaves,
there is exactly one isomorphism wherever there is one. The local
isomorphisms between the two pullbacks of $\xi$ therefore agree on overlaps
and glue. On $U\times_S U\times_S U$, the composite of two of them equals
the third because both have the same source and target. This is a descent
datum; effectiveness supplies the descended object. The same argument
descends its unique compatible isomorphisms. For a finite Galois extension,
the components of $\operatorname{Spec}L\times_k\operatorname{Spec}L$ are
indexed by the Galois group, so the assumed conjugacy isomorphisms are
exactly the required local isomorphisms.
\end{proof}

The hypotheses on local existence and effectiveness cannot be inferred
from absence of automorphisms alone. The heading does not specify a
category of structures, a topology, or a candidate ring; it cannot assert
existence of a least ring of definition for every conceivable structure.
The proposition isolates the precise descent obstruction which rigidity
does remove.

\subsection{(b) Hilbert and Picard without projectivity; the dual abelian scheme}

\begin{theorem}[Known: the nonprojective representability statements]
\label{geo-Kb-HilbPic}
Let $X\to S$ be a morphism of algebraic spaces.
If it is separated and finitely presented, the functor of closed subspaces
of $X_T$ proper, flat, and finitely presented over $T$ is an algebraic space,
locally of finite presentation over $S$.
If $X\to S$ is proper, flat, and finitely presented, its Picard stack is an
algebraic stack. If, in addition,
$\mathcal O_T\xrightarrow{\sim}(f_T)_*\mathcal O_{X_T}$ for every $T\to S$,
the fppf relative Picard sheaf is an algebraic space.
No projective embedding occurs in these assertions.
\end{theorem}
\begin{proof}[Derivation from opened results]
These are respectively Stacks Project, Proposition 99.9.4, tag 0D01;
Proposition 99.10.2, tag 0D04; and Proposition 99.11.8, tag 0D2C
\cite{geoStacksHilbert,geoStacksPicStack,geoStacksPicSpace}.
The functor in the first assertion requires proper support, not that $X$
itself be proper. The last assertion concerns the sheafification of
$T\mapsto\operatorname{Pic}(X_T)/\operatorname{Pic}(T)$; it does not
identify the groupoid of line bundles with a set-valued functor.
\end{proof}

\begin{theorem}[Known: duality without a projectivity hypothesis]
\label{geo-Kb-dual}
Let $A\to S$ be an abelian scheme, meaning a proper smooth group scheme
with geometrically connected fibres. Then $\operatorname{Pic}^0_{A/S}$ is
represented by an abelian scheme $A^\vee$, formation commutes with base
change, and there is a Poincar\'e bundle on $A\times_S A^\vee$ rigidified
on the two zero sections. Its associated biduality morphism
$A\to A^{\vee\vee}$ is an isomorphism.
\end{theorem}
\begin{proof}[Exact reduction and source audit]
The complete statement and reduction are printed in
\cite[Theorem 2.3.7, p.~15]{RobertAV2023}; the scheme-representability step
is Raynaud's theorem that an abelian algebraic space is a scheme, stated
there as Theorem 2.3.2, p.~12. The reduction starts with the Picard algebraic
space of Theorem~\ref{geo-Kb-HilbPic}. Its identity component is smooth and
proper with connected fibres; Raynaud's theorem makes it a scheme.
Rigidification along the zero section removes the pullback ambiguity in
the universal line bundle. Biduality is then checked on geometric fibres.
The source's sentence ``an abelian variety has no torsion'' must be read
as torsion-freeness of its \emph{N\'eron--Severi group}, which identifies
$\operatorname{Pic}^{\tau}$ with $\operatorname{Pic}^0$; abelian varieties
themselves certainly have torsion points.

For historical precision, the opened Orsay seminar
\cite[IX, Theorem 2.1; XI, 3.5]{RaynaudOrsay1967} proves formal smoothness
and proves the abelian-scheme conclusion \emph{conditional on}
representability by a scheme. It is not an independent proof of the last
representability step. The original Raynaud theorem remains a primary-text
retrieval gate in this run; the mathematical statement above has an opened,
precisely located modern source, not an inferred attribution from the
manuscript's marginal ``OK. par Raynaud''.
\end{proof}

\subsection{(c) Quotients without flatness}

\begin{theorem}[Proved correction and counterexample]
\label{geo-Kc}
Let $G\subseteq H=\mathrm{GL}_{n,S}$ be a finitely presented closed group
subscheme. Its fppf orbit sheaf $H/G$ is an algebraic space if and only if
$G\to S$ is flat.
In particular, arbitrary removal of flatness is false, even for a finite
commutative subgroup of $\mathrm{GL}_2$ over an Artinian base.
\end{theorem}
\begin{proof}
If $G$ is flat, the free-action quotient theorem
\cite[Lemma 80.11.7, tag 06PH]{geoStacksFreeQuotient} applies.
Conversely suppose $Q=H/G$ is an algebraic space. The smooth group $H$
acts on $Q$, so $[Q/H]$ is an algebraic stack. There is an equivalence
$[Q/H]\simeq BG$: a $G$-torsor gives its extended $H$-torsor and its
equivariant map to $H/G$; conversely the fibre over the identity coset of
such a map gives a $G$-torsor. The constructions are inverse after an fppf
cover on which the torsors and the coset lift are trivial, and therefore
globally by descent. This uses the defining sheaf-surjectivity of
$H\to H/G$, not flatness of that morphism. By
\cite[Lemma 97.18.3, tag 06PL]{geoStacksBG}, algebraicity of $BG$ forces
$G\to S$ flat and locally of finite presentation.

For the finite example set
\[
 A=\mathbf F_2[\epsilon]/(\epsilon^2),\qquad
 B=A[x]/(\epsilon x,x^2),\qquad G=\operatorname{Spec}B.
\]
Declare $x$ primitive: $\Delta(x)=x\otimes1+1\otimes x$,
$\varepsilon(x)=0$, and $\iota(x)=x$. These maps preserve the relations:
in characteristic two the square of the displayed sum is zero, and
$\epsilon\Delta(x)=0$. Coassociativity, the counit law, and the antipode
law follow directly on the generators. Thus $G$ is a finite commutative
group scheme, embedded in $\mathbf G_a$ and hence in $\mathrm{GL}_2$ by
$x\mapsto\left(\begin{smallmatrix}1&x\\0&1\end{smallmatrix}\right)$.
The $\mathbf F_2$-basis of $B$ is $1,\epsilon,x$. A finite flat module
over the local ring $A$ is free and has even $\mathbf F_2$-dimension;
$B$ has dimension three. It is not flat. Consequently its orbit quotient
$\mathrm{GL}_{2,A}/G$ is not an algebraic space. The finite Hopf-algebra
relations and multiplication are checked exhaustively by
\texttt{code/geometry\_finite\_checks.py}.
\end{proof}

\subsection{(d) A free action of $\mathrm{GL}_n$}

\begin{corollary}[Known, with the freeness condition made precise]
\label{geo-Kd}
If $\mathrm{GL}_{n,S}$ acts on an algebraic space $X$ with trivial
stabilizer functors, then $X/\mathrm{GL}_n$ is an algebraic space and
$X\to X/\mathrm{GL}_n$ is an fppf $\mathrm{GL}_n$-torsor.
\end{corollary}
\begin{proof}
The group is smooth, flat, and finitely presented over $S$; apply the
free-action theorem cited above. In its proof, the relation
$\mathrm{GL}_n\times_S X\rightrightarrows X$ has flat finitely presented
projections and maps monomorphically to $X\times_S X$ by freeness. The
quotient theorem for such equivalence relations gives exactly the asserted
algebraic space and torsor.
\end{proof}

``Sans point fixe'' must here mean free, not merely absence of a point
fixed by the entire group. For example $\mathbf G_m$ acting on itself by
$g\cdot x=g^2x$ over $\mathbf Q$ has no globally fixed point but has
stabilizer $\mu_2$. The corollary asserts an algebraic-space quotient;
representability by a scheme is an additional issue.

\subsection{(e) Moduli of arbitrary varieties, especially abelian schemes}

The words ``plus ou moins quelconques'' specify no class of varieties,
equivalence, topology, or finiteness condition. There is therefore no single
representability proposition to prove. Two precise broad interpretations
already fail.

\begin{example}[Refuted: the isomorphism-class presheaf is a moduli space]
\label{geo-Ke-twist}
The elliptic curves over $\mathbf R$
$E:y^2=x^3+x+1$ and $E':y^2=x^3+x-1$ are isomorphic over $\mathbf C$
and not over $\mathbf R$ as pointed curves. An isomorphism of short
Weierstrass equations in characteristic zero requires a nonzero $u$ with
$u^4=1$ and $u^6=-1$; equivalently $u^2=-1$. Thus the isomorphism-class
presheaf is not even separated for the fppf topology and cannot be
represented by an algebraic space. Both discriminants equal $-496$.
The exact polynomial check is in \texttt{geometry\_finite\_checks.py}.
Keeping the groupoid of objects and isomorphisms is necessary.
\end{example}

\begin{proposition}[Proved: the unrestricted abelian moduli stack fails]
\label{geo-Ke-unpolarized}
For $g\geq2$, the stack of all unpolarized abelian schemes of relative
dimension $g$ over $\mathbf C$ is not an algebraic stack locally of finite
presentation over $\mathbf C$.
\end{proposition}
\begin{proof}
Suppose a smooth atlas $U$ exists and choose a complex point $u$ over an
abelian variety $A$. By Theorem~\ref{geo-Ki}, deformations of $A$ are
unobstructed. Smoothness of the atlas then makes $U$ formally smooth at
$u$: first lift the family across a small Artinian extension, then use the
smooth lifting property of $U$ over the stack. Because $U$ is locally of
finite presentation over $\mathbf C$, it is smooth near $u$. Shrink to an
integral normal neighbourhood. The abelian scheme over this neighbourhood
is projective by \cite[\S3(C), p.~258]{Murre1965}, so carries an ample
line bundle and hence a polarization $\lambda$.

The differential from $T_uU$ to the marked deformation space of $A$ is
surjective because the atlas is smooth. It factors through deformations
preserving $\lambda$, since the whole family carries that polarization.
The unrestricted tangent dimension is $g^2$, while the polarized tangent
dimension in characteristic zero is $g(g+1)/2$.
These dimensions are given in \cite[\S2.2, Theorem 2.2.1;
\S2.4, Theorem 2.4.1 and its proof]{Oort1971}.
Equivalently, a polarization identifies the tangent space with matrices
and its first-order lifting condition is symmetry. Since
$g(g+1)/2<g^2$, the required surjectivity is impossible.
\end{proof}

This proposition explains why a chosen polarization or an appropriate
bounded moduli problem is a material additional hypothesis. It does not
assert that all moduli problems without a polarization fail, nor confuse
the smooth formal deformation space with an algebraic atlas.

\subsection{(f) Principal bundles, commutative and noncommutative}

\begin{theorem}[Proved classification; known algebraicity under precise hypotheses]
\label{geo-Kf}
For an affine group scheme $G$ on $Y$, principal $G$-bundles on $X$ in the
fppf topology are classified by $H^1_{\mathrm{fppf}}(X,G_X)$, interpreted
as the set of torsor isomorphism classes. This set is an abelian group if
$G$ is commutative; in general it is a pointed set, and the classifying
groupoid retains the automorphisms.
If $X\to Y$ is proper, flat, and finitely presented and $G\to Y$ is
affine, flat, and finitely presented, the relative stack of these bundles
is algebraic, locally of finite presentation over $Y$, with affine diagonal.
\end{theorem}
\begin{proof}
Trivialize a torsor on an fppf covering $\{U_i\to X\}$. The changes of
trivialization give $g_{ij}\in G(U_i\times_XU_j)$ satisfying
$g_{ij}g_{jk}=g_{ik}$. Changing trivializations by $h_i$ replaces them by
$h_i^{-1}g_{ij}h_j$. Conversely these equations are precisely effective
descent data for the trivial affine $G$-schemes and their right actions.
Refining coverings identifies presentations of the same torsor. This
constructs the stated classification, including morphisms. For a
commutative group, contracted product of torsors gives addition, the
trivial torsor gives zero, and inversion gives inverses; the cocycle law
makes these operations independent of choices. Without commutativity,
pointwise multiplication of two cocycles need not be a cocycle.

For the second assertion, $BG$ is an algebraic stack by flatness and finite
presentation of $G$. Its diagonal is affine, and it is locally of finite
presentation and quasi-separated over $Y$. The desired stack is
$\operatorname{Hom}_Y(X,BG)$. Apply Hall--Rydh,
\cite[Theorem 1.2, p.~1634]{HallRydh2019}, which assumes exactly a proper
flat finitely presented source and a locally finitely presented
quasi-separated target with affine stabilizers. Its part (iii) preserves
affineness of the diagonal. Projectivity of $X$ is unnecessary.
\end{proof}

\subsection{(g) The finite-group case}

\begin{proposition}[Known classification and proved inseparable correction]
\label{geo-Kg}
For a connected scheme $X$ with geometric point $\bar x$ and a finite
constant group $G$, its $G$-torsors are classified by continuous
homomorphisms $\pi_1^{\mathrm{\acute et}}(X,\bar x)\to G$ modulo
conjugacy. A chosen point of the torsor above $\bar x$ removes conjugacy.
For a ring $R$ of characteristic $p$ one instead has
\[
 H^1_{\mathrm{fppf}}(\operatorname{Spec}R,\alpha_p)=R/R^p.
\]
Thus ``finite'' cannot silently be replaced by ``finite \'etale''.
\end{proposition}
\begin{proof}
The equivalence of finite \'etale covers and finite continuous fundamental
group sets is \cite[Theorem 58.6.2, tag 0BND]{geoStacksPiOne}.
A torsor fibre is a simply transitive right $G$-set. Choosing one point
identifies its commuting fundamental-group action with left multiplication
by a homomorphism to $G$; changing the point conjugates it.

For the second assertion, the fppf exact sequence
$0\to\alpha_p\to\mathbf G_a\xrightarrow{F}\mathbf G_a\to0$
is exact on the right because $z^p-f$ has a root after the faithfully flat
finite free cover $R[z]/(z^p-f)$. The vanishing of the first cohomology of
the quasi-coherent additive sheaf on an affine scheme gives the displayed
cokernel. Concretely the torsor for $f$ is $z^p=f$, and changing $z$ by
$b\in R$ changes $f$ by $b^p$. For a perfect field $k$ and $R=k[t]$, the
classes of $t^m$ with $p\nmid m$ are linearly independent over $k$ in this
quotient. Finiteness of the group scheme does not make the set of bundles
finite.
\end{proof}

\subsection{(h) $X=\mathbf P^1_Y$, $G=\mathrm{GL}_n$}

\begin{theorem}[Proved: field classification and exact relative replacement]
\label{geo-Kh}
Every rank $n$ vector bundle on $\mathbf P^1_k$, for any field $k$, is
$\bigoplus_{i=1}^n\mathcal O(a_i)$, with its integers determined up to
permutation. Over an arbitrary affine base $Y=\operatorname{Spec}R$,
the complete classification is instead the groupoid of triples
\[
 (P_0,P_\infty,\phi),\qquad
 \phi:P_0\otimes_{R[t]}R[t,t^{-1}]
 \xrightarrow{\sim}P_\infty\otimes_{R[t^{-1}]}R[t,t^{-1}],
\]
where the two modules are finitely generated projective of rank $n$;
morphisms are pairs of module isomorphisms intertwining $\phi$.
For general $Y$, these data are descended along an affine open covering.
Principal $\mathrm{GL}_n$-bundles have exactly this classification.
\end{theorem}
\begin{proof}
For the field assertion choose the greatest integer $a$ for which
$H^0(E(-a))\ne0$. Such integers exist and are bounded above, for instance
by trivializing on the two standard affine lines and bounding the powers
appearing in a Laurent transition matrix and its inverse. A nonzero
section gives $\mathcal O(a)\to E$. Saturating this line subbundle would
raise its degree if the section had a zero, contradicting maximality of
$a$; thus the quotient $Q$ is locally free. By induction
$Q=\bigoplus_i\mathcal O(b_i)$. In the exact sequence twisted by $-a-1$,
$H^0(E(-a-1))=0$ and $H^1(\mathcal O(-1))=0$, so
$H^0(Q(-a-1))=0$. Hence $b_i\leq a$ for every $i$.
The extension splits because
$\operatorname{Ext}^1(\mathcal O(b_i),\mathcal O(a))
=H^1(\mathcal O(a-b_i))=0$. The dimensions of $H^0(E(m))$ for all $m$
recover the multiset of degrees, proving uniqueness.

For the relative assertion, the standard open cover has coordinate rings
$R[t]$, $R[t^{-1}]$, and overlap $R[t,t^{-1}]$. Vector bundles on an affine
scheme are finite projective modules, and effective Zariski descent gives
exactly the triples and their morphisms. Their chart modules need not be
free, so a matrix double-coset description without that qualification would
omit objects. Taking frames and forming the associated rank $n$ vector
bundle are inverse equivalences between $\mathrm{GL}_n$-torsors and
vector bundles.
\end{proof}

\begin{example}[Refuted: a relative bundle has constant splitting type]
\label{geo-Kh-jump}
On $\mathbf P^1_k\times\mathbf A^1_k$ take the extension
$0\to\mathcal O(-1)\to E\to\mathcal O(1)\to0$
whose class is the coordinate $s$ in
$H^1(\mathbf P^1,\mathcal O(-2))\otimes_k k[s]\simeq k[s]$.
At $s=0$ it is $\mathcal O(-1)\oplus\mathcal O(1)$.
At a nonzero field-valued point the boundary map after twisting by $-1$
is multiplication by $s$ from $H^0(\mathcal O)$ to
$H^1(\mathcal O(-2))$. Consequently $H^0(E(-1))=0$, and the field
classification with degree zero forces $E\simeq\mathcal O^2$.
This is the precise obstruction behind treating (h) as just a list of
integers independent of the base. The extension is locally free because
it is an extension of locally free sheaves.
\end{example}

\subsection{(i) Are the formal moduli of an abelian variety smooth?}

\begin{theorem}[Known, with all characteristics and the base specified]
\label{geo-Ki}
Let $A_0/k$ be an abelian variety of dimension $g$, and let $W$ be a
complete local noetherian ring with residue field $k$. On local Artinian
$W$-algebras with residue field identified with $k$, the functor of abelian
schemes equipped with an identification of their special fibre with $A_0$
is prorepresented by
\[
 W[[t_{11},t_{12},\ldots,t_{gg}]].
\]
In particular, it is formally smooth over $W$, of relative tangent
dimension $g^2$. No polarization is part of this deformation problem.
\end{theorem}
\begin{proof}[Derivation from the exact primary theorem]
This is Oort's Theorem 2.2.1, printed p.~273, with the coefficient category
defined in \S2.1, p.~269 \cite{Oort1971}. Its proof on pp.~273--280
handles characteristic two explicitly by an automorphism of $A_0\times A_0$;
the shorter argument using $[-1]$ alone would leave that characteristic
unproved. Formal smoothness now follows directly: across a surjection of
Artinian coefficient algebras, lift the images of the finitely many
parameters. The tangent dimension is also
$\dim H^1(A_0,T_{A_0})=\dim(H^1(A_0,\mathcal O_{A_0})
\otimes_k\operatorname{Lie}A_0)=g^2$.
\end{proof}

Here ``simples'' is interpreted in the historical sense of smoothness.
This formal conclusion does not assert algebraization of every formal
family or existence of an unrestricted algebraic moduli space; the
distinction is forced by Proposition~\ref{geo-Ke-unpolarized}.

\subsection{(j) The theorem of the square without projectivity}

\begin{theorem}[Proved from the nonprojective cube theorem]
\label{geo-Kj}
Let $p:A\to S$ be an abelian scheme with zero section $e$, let $L$ be an
invertible sheaf, and let $x,y\in A(S)$. Then
\begin{align*}
 t_{x+y}^*L\otimes(t_x^*L)^{-1}\otimes(t_y^*L)^{-1}\otimes L
 \simeq p^*\bigl((x+y)^*L\otimes(x^*L)^{-1}
 \otimes(y^*L)^{-1}\otimes e^*L\bigr).
\end{align*}
Thus $x\mapsto[t_x^*L\otimes L^{-1}]$ is a homomorphism into the relative
Picard functor. Neither this assertion nor its proof requires projectivity.
\end{theorem}
\begin{proof}
The relative theorem of the cube is
\cite[VII, Theorem 3.1, pp.~VII-6--VII-8]{RaynaudOrsay1967}. It applies to
proper flat factors with sections and with universally
$f_*\mathcal O=\mathcal O$; these hypotheses hold for an abelian scheme.
On $A^3$ form
\[
 m_{123}^*L\otimes(m_{12}^*L)^{-1}\otimes(m_{13}^*L)^{-1}
 \otimes(m_{23}^*L)^{-1}\otimes p_1^*L\otimes p_2^*L
 \otimes p_3^*L\otimes(p e)^*L^{-1}.
\]
Here $m_I$ adds the coordinates in $I$, and the last factor denotes
pullback of $e^*L$ from $S$. On each coordinate hyperplane with one
coordinate zero, factors cancel canonically. The cube theorem makes this
rigidified line bundle trivial. Pull it back by
$a\mapsto(a,x(p(a)),y(p(a)))$ and rearrange its eight factors; the
displayed formula results. Its right side is pulled back from the base,
so its class in the relative Picard functor is zero. This proves the
homomorphism law after every base change.
\end{proof}

\begin{example}[Refuted: suppressing the base line bundle]
\label{geo-Kj-base}
Let $S$ be a complex elliptic curve $E$, take $A=E\times S$, let
$L=\operatorname{pr}_1^*\mathcal O_E(0)$, and take $x=y$ to be the
diagonal section. The base line bundle in Theorem~\ref{geo-Kj} has degree
$4-1-1=2$, and is nontrivial. Its pullback to $A$ is nontrivial because
the zero section pulls it back to itself. Thus the absolute equality in
$\operatorname{Pic}(A)$ without that base factor is false even for a
projective constant family. Over a field the factor is a one-dimensional
vector space and becomes trivial as a line bundle, which explains the
usual field-valued formulation.
\end{example}
